\section*{8 Differential Calculus in Several Variables}
\begin{enumerate}
    \setcounter{enumi}{4}
    \renewcommand{\theenumi}{\alph{enumi}}
    \item Let
    \[
    D^\alpha f(\mathbf{x}) := \frac{\partial^{|\alpha|} f}{(\partial x^1)^{\alpha_1} \ldots (\partial x^m)^{\alpha_m}}(\mathbf{x}).
    \]
    Show that if $f \in C^{(k)}(G; \mathbb{R})$, then the equality
    \[
    \sum_{i_1 + \ldots + i_m = k} \frac{k!}{i_1! \ldots i_m!} f^{(i_1, \ldots, i_m)} h^{i_1} \ldots h^{i_m} = \sum_{|\mathbf{\alpha}|=k} \frac{k!}{\mathbf{\alpha}!} D^\mathbf{\alpha} f(\mathbf{x}) h^\mathbf{\alpha},
    \]
    where $\mathbf{h} = (h^1, \ldots, h^m)$, holds at any point $\mathbf{x} \in G$.
    \item Verify that in multi-index notation Taylor's theorem with the Lagrange form
    of the remainder, for example, can be written as
    \[
    f(\mathbf{x} + \mathbf{h}) = \sum_{|\mathbf{\alpha}|=0}^{n-1} \frac{1}{\mathbf{\alpha}!} D^\mathbf{\alpha} f(\mathbf{x}) \mathbf{h}^\mathbf{\alpha} + \sum_{|\mathbf{\alpha}|=n} \frac{1}{\mathbf{\alpha}!} D^\mathbf{\alpha} f(\mathbf{x} + \theta \mathbf{h}) \mathbf{h}^\mathbf{\alpha}.
    \]
    \item Write Taylor's formula in multi-index notation with the integral form of the
    remainder (Theorem 4).
\end{enumerate}
\begin{enumerate}
    \setcounter{enumi}{5}
    \item a) Let $I^m = \{\mathbf{x} = (x^1, \ldots, x^m) \in \mathbb{R}^m \mid |x^i| \leq c^i, i = 1, \ldots, m\}$ be an $m$-dimensional closed interval and $I$ a closed interval $[a, b] \subset \mathbb{R}$. Show that if the function
    $f(\mathbf{x}, y) = f(x^1, \ldots, x^m, y)$ is defined and continuous on the set $I^m \times I$, then for any
    positive number $\varepsilon > 0$ there exists a number $\delta > 0$ such that $|f(\mathbf{x}, y_1) - f(\mathbf{x}, y_2)| <
    \varepsilon$ if $\mathbf{x} \in I^m, y_1, y_2 \in I$, and $|y_1 - y_2| < \delta$.
    \item b) Show that the function
    \[
    F(\mathbf{x}) = \int_{a}^{b} f(\mathbf{x}, y) dy
    \]
    is defined and continuous on the closed interval $I^m$.
    \item c) Show that if $f \in C(\mathbb{I}^m; \mathbb{R})$, then the function
    \[
    F(\mathbf{x}, t) = f(t\mathbf{x})
    \]
    is defined and continuous on $I^m \times I^1$, where $I^1 = \{t \in \mathbb{R} \mid |t| \leq 1\}$.
    \item d) Prove Hadamard's lemma:
    If $f \in C^{(1)}(\mathbb{I}^m; \mathbb{R})$ and $f(\mathbf{0}) = \mathbf{0}$, there exist functions $g_1, \ldots, g_m \in C(\mathbb{I}^m; \mathbb{R})$
    such that
    \[
    f(x^1, \ldots, x^m) = \sum_{i=1}^m x^i g_i(x^1, \ldots, x^m)
    \]
    in $I^m$, and in addition
    \[
    g_i(0) = \frac{\partial f}{\partial x^i}(0), \quad i = 1, \ldots, m.
    \]
\end{enumerate}